En una zona de l'espai hi ha un camp magnètic uniforme de valor $\vec{B} = 2,00\,\vec{k}\ \mathrm{mT}$. Un electró, un neutró i un protó hi entren per l'origen de coordenades a la mateixa velocitat $\vec{v} = 5,00\,\vec{\jmath}\ \mathrm{m\,s^{-1}}$.

\begin{apartat}{1}
Determineu el mòdul de la força que actua sobre cada partícula i indiqueu el tipus de moviment que fa cadascuna.
\end{apartat}

\begin{apartat}{1}
A continuació, situem paral·lelament a l'eix $Y$, a 3,00~mm de l'origen de coordenades, un fil infinit pel qual circula un corrent $I$. Determineu el valor del corrent del fil que fa que el protó segueixi una trajectòria rectilínia. Considereu que el mòdul del camp magnètic creat per aquest fil infinit és: $B = \dfrac{\mu_0 I}{2\pi R}$, en què $\mu_0 = 4\pi\times 10^{-7}\ \mathrm{N\,A^{-2}}$ i $R$ és la distància al fil conductor.
\end{apartat}

\dades{%
  Càrrega elèctrica del protó, $q_\text{protó} = 1,60\times 10^{-19}\ \mathrm{C}$.\\
  Càrrega elèctrica de l'electró, $q_\text{electró} = -q_\text{protó}$.}
